\documentclass[12pt]{article}
\usepackage[utf8]{inputenc}
\usepackage[T1]{fontenc}
\usepackage{hyperref}
\title{Kujarge Language Resources for K-12 Curriculum Development}
\author{Compiled by Claude Sonnet 5 (claude-sonnet-5) for Andr\'as Kornai}
\date{}
\begin{document}
\maketitle

\section{Introduction}
Kujarg\'e (Kujarke) is an endangered, very poorly documented language of eastern Chad (Ouadda\"i region, near Jebel Mirra), with related villages in Darfur, Sudan. Its classification is unresolved. The only speaker estimate found is about 1,000, dating from 1983. Sources disagree on whether the villages are in Chad or Sudan, and war in Darfur means the current situation is unknown.

\section{Primary Educational Resources}
I found no primers, textbooks or school programmes for Kujarg\'e.

\section{Linguistic and Reference Materials}
\begin{itemize}
\item Kujarg\'e language (Wikipedia): \url{https://en.wikipedia.org/wiki/Kujarg%C3%A9_language}
\item Kujarke people (Wikipedia): \url{https://en.wikipedia.org/wiki/Kujarke_people}
\item Endangered Languages Project entry: \url{https://endangeredlanguages.com/lang/575?hl=en}
\item V.\ Bla\v{z}ek on Kujarke as a possible independent East Chadic group (Folia Orientalia 52): \url{https://journals.pan.pl//Content/109618/PDF/FoliaOrientalia%2052-15%209%20Blazek.pdf}
\item East Chadic B classification and description, progress report (Journal of West African Languages 40.1): \url{https://journalofwestafricanlanguages.org/downloads/download/116-volume-40-number-1/601-east-chadic-b-classification-and-description-progress-report}
\end{itemize}
The language is known mainly from a roughly 200-word list by Doornbos (1981).

\section{Religious and Community Resources}
\begin{itemize}
\item Joshua Project profile, Kujarge in Chad: \url{https://joshuaproject.net/index.php/people_groups/11266}
\end{itemize}

\section{Digital and Audio Resources}
No audio, video or software resources were found.

\section{Curriculum Development Context}
No educational materials exist, and the language is not known to be taught in schools. Basic documentation (word lists, texts, recordings, an orthography) would have to come first, in consultation with the community and with attention to the security situation. The language of schooling in the area would presumably be Arabic or French.
\end{document}